\section*{Notations}
\addcontentsline{toc}{section}{Notations}

\begin{center}
\renewcommand{\arraystretch}{1.25}
\begin{tabular}{>{\raggedright}p{0.3\textwidth} p{0.6\textwidth}}
\toprule
\textbf{Notation} & \textbf{Signification} \\
\midrule
$x\in\R^d$ & vecteur (en colonne) de $d$ nombres $x_1,\dots,x_d$ \\
$A\T$ & transposée de la matrice $A$ (lignes et colonnes échangées) \\
$\sum_{i=1}^n a_i$ & somme $a_1+a_2+\dots+a_n$ \\
$\bar y=\frac1n\sum_{i=1}^ny_i$ & moyenne des $n$ nombres $y_1,\dots,y_n$ \\
$\P(A)$ & probabilité de l'événement $A$ \\
$\E[X]$, $\Var(X)$ & espérance (moyenne théorique) et variance de $X$ \\
$\P(A\mid B)$, $\E[Y\mid X=x]$ & probabilité de $A$ sachant $B$ ; moyenne de $Y$ sachant $X=x$ \\
$\mathcal N(\mu,\sigma^2)$ & loi normale de moyenne $\mu$ et de variance $\sigma^2$ \\
$\Phi$ & fonction de répartition de la loi $\mathcal N(0,1)$ \\
$\argmin_c g(c)$ & la (ou les) valeur(s) de $c$ où $g$ atteint son minimum \\
\bottomrule
\end{tabular}
\end{center}

Les variables aléatoires sont notées en majuscules ($X$, $Y$), les valeurs observées en minuscules ($x$, $y$).
